\section{Noir et blanc}
Le \robot\ peut colorier ses dessins. 
Nous allons commencer par des coloriages en noir et blanc. La commande {\idt peindre} 
permet de laisser s'échapper la peinture noire par l'arrière du \robot. La peinture
s'étale sur la plus grande partie possible de l'écran délimitée par une ligne noire. 
Un exemple est donné à la figure \ref{fig:peinyy} ci-dessous. 
\begin{figure}[htbp]
 \includevspacePostScript[2.9cm]{PeintureYinYang}{0.95}{0.95}
 \caption{Étapes de la peinture d'un dessin}\label{fig:peinyy}
\end{figure}
On commence par dessiner le pourtour de la zone à peindre puis on place le \robot\ dans la zone à
peindre ou sur son bord 
de façon 
à ce que l'arrière du \robot\ par où s'échappe la peinture soit dans la zone à peindre. 
La peinture s'étale ligne par ligne au dessus,
à gauche, à droite et en dessous du point arrière du \robot. La zone couverte est la plus grande 
où se trouve 
l'arrière du 
\robot\ et qui est
délimitée par une ligne noire continue. Le \robot\ peut peindre sans se déplacer, que son crayon soit levé ou baissé.
\begin{exemple}[Yin Yang]{\rm
 Le programme ci-après dessine le symbole chinois du Yin Yang~:\hfill
 \pagebreak[3]
 \vbox to 0pt{\vspace*{25mm}\includePostScript{YinYang}{0.95}{0.95}}
 \input YinYang.tex
\noindent $\Box$ }\end{exemple}

Une erreur fréquente consiste à laisser un trou dans le périmètre de la zone à peindre par
lequel la peinture fuira. Un trou d'un seul pixel suffit pour provoquer une fuite. En cas de fuite, on peut
arrêter la peinture en tapant sur l'une des touches
\begin{picture}(12,12)(0, 2) %touche d'echappement
  \put(0,0){\framebox(12,12)[lb]{}}
  \put(1.5,2){\makebox(0,0)[lb]{\tiny esc}}
  \end{picture}\hspace{0.3ex} ou 
  \begin{picture}(12,12)(0, 2) %touche §
  \put(0,0){\framebox(12,12)[lb]{}}
  \put(1.5,2){\makebox(0,0)[lb]{\tiny \S}}
  \end{picture}\hspace{0.3ex} ou
  \begin{picture}(12,12)(0, 2) %touche #
  \put(0,0){\framebox(12,12)[lb]{}}
  \put(1.5,2){\makebox(0,0)[lb]{\tiny \#}}
  \end{picture}\hspace{0.3ex}. Le \robot\ continue de dessiner.
  Taper au moins une seconde fois sur l'une de ces touches pour l'arrêter définitivement.

\section{Couleur}

Pour dessiner et peindre, le \robot\ dispose de crayons de couleurs blanc, jaune, bleu, violet, vert, 
rouge, indigo ou noir. La commande {\idt cc} permet de choisir la \underline{c}ouleur du 
\underline{c}rayon, que l'on donne en paramètre. Par exemple après la commande {\idt cc(rouge)} 
tous les tracés et la peinture sont en rouge. La commande {\idt cc(noir)} permet de
revenir aux tracés et peinture noirs.

Sur un écran couleur, le \robot\ dessine et peint avec les couleurs de la première
ligne ci-dessous.

Quand on ne dispose pas d'un écran couleur mais d'un écran à plusieurs niveaux de gris
dont chaque pixel peut être plus ou moins foncé, les
couleurs sont représentées par des gris, dans l'ordre blanc, 
jaune, bleu, violet, vert,  rouge, indigo et noir du plus clair au plus foncé.

Quand on dispose seulement d'un écran noir et blanc, les couleurs sont représentées par des
grisés constitués par des motifs de points blancs et de points noirs juxtaposés. L'ordre est
à nouveau blanc, 
jaune, bleu, violet, vert,  rouge, indigo et noir du plus clair au plus foncé. Plus les grisés sont
foncés plus les petits points utilisés pour les représenter sur l'écran sont denses comme
on le voit sur la deuxième ligne ci-dessous~:\\[1.8cm]
\includePostScript{CouleursCouleur}{0.95}{0.95}\ \\[1.8cm]
\includePostScript{Couleurs}{0.95}{0.95}\ \\

Quand on peint sur un écran noir et blanc, les bords du motif doivent être tracés en noir (car
dans une autre couleur le pourtour serait tracé en pointillés ce qui provoquerait
des fuites de peinture). 

\begin{exemple}[Fleur en couleur]{\rm
 Le programme ci-dessous permet de peindre la fleur ci-contre en couleur~:\\[5cm]
 \includePostScript{FLeurStyliseeCouleur}{0.95}{0.95}%
 \vspace*{-5cm}
 \input FLeurStyliseeCouleur.tex
\noindent $\Box$ }\end{exemple}

\noindent \rule{137mm}{0.1mm} %---------------------------------------------
\begin{exercices}{Écrire des programmes \pascal\ pour dessiner et colorier en noir et blanc~:

\tiret les pavages du carré de la figure \ref{fig:pavac} ci-dessous$\:$;

\tiret les mosaïques de la salle des Abencérages (XIV\/$^{\rm\grave{e}me}$ siècle) de l'Ah\-lam\-bra, 
la résidence des rois maures à Grenade en Andalousie telles qu'elles ont été relevées
par le peintre Escher en 1936 \cite{escher1} (figure \ref{fig:mosAlham} de la page \pageref{fig:mosAlham}).

\noindent $\Box$}\end{exercices}
\begin{figure}[htbp]
 \includevspacePostScript[4.3cm]{Pavages.DuCarre}{0.95}{0.95}
 \caption{Pavages du carré}\label{fig:pavac}
\end{figure}
\begin{figure}[htbp]
 \includevspacePostScript[4.6cm]{MosaiqueAlhambra123}{0.95}{0.95}
 \caption{Mosaïques de l'Alhambra}\label{fig:mosAlham}
\end{figure}
\begin{exercices}{Écrire des programmes \pascal\ pour dessiner et colorier en cou\-leur~:

\tiret la balle et le carrelage en trois couleurs de la figure \ref{fig:balpav} (page \pageref{fig:balpav})$\:$;

\tiret la mosaïque de l'Ah\-lam\-bra \cite{escher1} de la figure \ref{fig:mosai4} 
        (page \pageref{fig:mosai4})$\:$;

\tiret le pavage de carrés disjoints de la figure \ref{fig:pdcc} (page \pageref{fig:pdcc})$\:$;

\tiret les pavages des figures 
\ref{fig:PVdV} (page \pageref{fig:PVdV}),
\ref{fig:PVdV1l} (page \pageref{fig:PVdV1l}),
\ref{fig:PVdV2l} (page \pageref{fig:PVdV2l})
et
\ref{fig:PVdV4l} (page \pageref{fig:PVdV4l})
utilisant le motif ornemental de N. J. van de Vecht comme il est
peint à la figure \ref{fig:mVdV} (page \pageref{fig:mVdV}).

\noindent $\Box$}\end{exercices}
\begin{figure}[htbp]
\includevspacePostScript[4.4cm]{BallePavageCouleur}{0.95}{0.95}
 \caption{Balle, pavage en trois couleurs}\label{fig:balpav}
\end{figure}
\vfill
\begin{figure}[htbp]
 \includevspacePostScript[5.6cm]{MosaiqueAlhambra4}{0.95}{0.95}
 \caption{Mosaïque de l'Alhambra en couleur}\label{fig:mosai4}
\end{figure}
\begin{figure}[htbp]
\vspace*{\textwidth}%
\vspace*{25mm}%
\includevspacePostScript[8.8cm]{PavageDis.Car.Cou}{0.95}{0.95}%
\vspace*{-\textwidth}%
\vspace*{-25mm}%
\caption{Pavage de carrés disjoints en couleur}\label{fig:pdcc}
\end{figure}
\begin{exercice}{Quel est le motif utilisé pour réaliser le pavage noir et blanc de la figure \ref{fig:pavhexond}
(page \pageref{fig:pavhexond})$\:$?
\noindent $\Box$}\end{exercice}
\begin{figure}[htbp]
 \includevspacePostScript[13.7cm]{PavageHexagonalOndule}{0.95}{0.95}
 \caption{Pavage du plan en noir et blanc}\label{fig:pavhexond}
\end{figure}
%%%%%%%%%%%%%%%%%%%%%%%%%%%%%%%%%%%%%%%%%%%%%%%%%%%%%%%%
\begin{corrige}{
Nous donnons un programme, les autres se trouvent sur la disquette  d'accompagnement du livre~:
\input MosaiqueAlhambra3.tex
\noindent $\Box$}\end{corrige}
\begin{corrige}{
Le motif du pavage du plan en noir et blanc de la figure \ref{fig:pavhexond}
(page \pageref{fig:pavhexond}) est donné à la figure \ref{fig:motpavhexond}
(page \pageref{fig:motpavhexond}).
\noindent $\Box$}\end{corrige}
\begin{figure}[htbp]
 \includevspacePostScript[7.4cm]{MotifDuPavageHexagonalOndule}{0.95}{0.95}
 \caption{Motif du pavage de la figure \protect{\ref{fig:pavhexond}}}\label{fig:motpavhexond}
\end{figure}
\newpage
\begin{corrige}{
Nous donnons un programme, les autres se trouvent sur la disquette  d'accompagnement du livre~:
\input BalleEnCouleur-2.tex
\noindent $\Box$}\end{corrige}

\vfill

%%%%%%%%%%%%%%%%%%%%%%%%%%%%%%%%%%%%%%%%%%%%%%%%%%%%%%%%%%%%%%%%%%%%%%%%%%%
\noindent \framebox[\textwidth]{
\begin{minipage}{\myparboxwidth}
\setlength{\parindent}{2em}
\small \sf
\vspace{\abovedisplayskip}
\section*{Addenda}

On peut imaginer que pour peindre, le \robot\ place un germe sur l'écran, en arrière
de son point d'arrêt. Si un germe est placé sur un pixel de l'écran déjà colorié alors 
ce germe disparaît. En particulier si le germe initial placé par le \robot\ est sur un pixel non blanc,
la peinture s'arrête immédiatement.
Si le germe est placé sur un pixel blanc de l'écran, alors le germe disparaît, le pixel blanc sur
lequel il se trouvait est peint de la couleur du
crayon et quatre nouveaux germes sont placés à l'est, à l'ouest, au nord et au sud du pixel qui
vient d'être peint. La peinture est terminée quand tous les germes ont disparu.

Cet algorithme de peinture est donc basé sur la représentation physique du dessin sur l'écran et
non pas sur la représentation logique qui serait celle obtenue en mémorisant toutes les commandes
exécutées par le \robot. De ce fait il se produit parfois des fuites logiquement inexplicables mais
qui sont dues à la différence des représentations logiques et physiques.

Prenons l'exemple d'un cercle. Mathématiquement un cercle est une ligne continue cons\-ti\-tuée
par un ensemble infini de points sans dimension. Sur l'écran d'un ordinateur le cercle sera
représenté par une ligne discontinue cons\-ti\-tuée par un ensemble fini de "points" de la dimension
d'un pixel. Les notions mathématique et informatique de cercle sont donc complètement
différentes.

 L'approximation est acceptable quand la taille des pixels est très faible par rapport
au diamètre du cercle de la même façon que le compas permet de tracer des cercles qui
donnent une bonne image de l'objet mathématique. 
Cependant pour des cercles très petits, on
arrive à quelques aberrations visibles sur la figure ci-après où les cercles représentés
dans le petit cadre en haut à gauche de la figure ont été agrandis  huit fois.
On notera en particulier que les très petits cercles sont carrés$\:$!
\vspace{\abovedisplayskip}
\end{minipage}
}
%%%%%%%%%%%%%%%%%%%%%%%%%%%%%%%%%%%%%%%%%%%%%%%%%%%%%%%%%%%%%%%%%%%%%%%%%%%%
%%%%%%%%%%%%%%%%%%%%%%%%%%%%%%%%%%%%%%%%%%%%%%%%%%%%%%%%%%%%%%%%%%%%%%%%%%%
\noindent \framebox[\textwidth]{
\begin{minipage}{\myparboxwidth}
\setlength{\parindent}{2em}
\small \sf
\vspace{\abovedisplayskip}

\includehvspacePostScript[9.8cm,5.2cm]{LoupeCercles}{0.75}{0.75}

\vspace*{5mm}
Sur les écrans ayant des pixels rectangulaires, les cercles ont des
formes d'ellipses.  Le programme doit donc corriger la déformation en
dessinant logiquement des ellipses pour obtenir des cercles sur
l'écran.  Les cercles sont fermés en ce sens que tout pixel peint sur
le cercle a au moins un pixel voisin peint dans l'une des huit
directions de la rose des vents.  Ce n'est pas le cas pour les
ellipses très allongées que l'on peut voir sur la figure ci-dessous~:
\vspace*{5mm}

\includehvspacePostScript[9.7cm,4.0cm]{LoupeEllipses}{0.75}{0.75}

\vspace*{5mm}
Les trous dans le périmètre des ellipses allongées sont des points de fuite de la peinture
quand on veut peindre l'intérieur de l'ellipse. Pour éviter ces fuites il faudrait utiliser des
algorithmes de coloriage plus complexes basés sur la structure logique et non pas physique du 
dessin.
\vspace{\abovedisplayskip}
\end{minipage}
}
%%%%%%%%%%%%%%%%%%%%%%%%%%%%%%%%%%%%%%%%%%%%%%%%%%%%%%%%%%%%%%%%%%%%%%%%%%%%
